\section{Conclusion}

We asked whether a small model can keep improving on the job, overnight, from the reactions of the people it works for, with no one checking the result. Over 127 lives, the answer turned on the judge: lives kept their gains when failed attempts stayed out of their lessons, and a few wrong lessons could undo them. The model cannot be its own judge; it learns and then decays. A separate, frozen reader of the user can stand in for the true outcome for a while, if it has discernment: if it waits for a clear verdict on the result instead of taking politeness for approval. Even so, by day 60 ours was no better than self-judging. Discernment has limits. It cannot recover what the user did not see, and where the user cannot check the result, learning from users teaches failure; and even correct lessons can fade when scarce, and may narrow the model's skills when lopsided. In our experiments, learning from users held up only where the judge kept failures out and the user could see the outcome. The data and code to test this further are public.
